% =============================================================
\section{Discussion}\label{sec:discussion}
% =============================================================

Genetic search over
waveforms~\cite{grill2015model,wongsarnpigoon2010energy} returns no
gradient direction for a candidate that is close to but short of an
activation target; surrogate networks trained on NEURON
outputs~\cite{hussain2024efficient} introduce an approximation error
that must be re-incurred whenever the underlying biophysical model
changes. Because the pipeline is JAX-traceable end-to-end, gradients
$\partial \ell / \partial w(t)$ of any differentiable scalar objective
$\ell$ with respect to the input waveform are returned in a single
backward pass and avoid both costs. As a
sharper receipt than the position gradient of \cref{sec:res-grad}, an
Adam loop on $\mathrm{MSE}(\hat{v}, v_\mathrm{target})$ on the full
BBP Pyr morphology recovers \texttt{gNaTs2T} from a $-15\,\%$
perturbation to within $+0.13\,\%$ in nine steps
(\texttt{fit\_demo}); the point of this fit is end-to-end gradient
flow on a realistic cell, not optimiser performance. A selectivity
Pareto sweep against a non-target cell, or a charge-minimising waveform
under a spike-elicitation constraint, follows the same template (a
differentiable spike-elicitation surrogate plus charge-balance and
amplitude constraints).
